% AI-produced by GPT-6 Astra, OpenAI. Reconstructed 20 November 2026.
% Mathematical reconstruction of the pre-outage A150 proof and its
% Hausdorff consequence, previously audited by root and singularities_duality.
% The original Valette v7 PDF was reopened during this reconstruction.
\Needspace{10\baselineskip}
\subsection{A.150. A definable scalar clock for a spatially Lipschitz family}

\begin{saresult}{A.150}{Proved under stated hypotheses}{Scalar-clock regularization and its Hausdorff consequence}
Work in any fixed o-minimal expansion of the real field.
Let $Z\subset\mathbb R^n$ be compact and definable, and let
$F:[0,1]\times Z\to\mathbb R$ be continuous and definable.
Suppose that for one $L\ge0$,
\[
 |F(t,z)-F(t,z')|\le L\|z-z'\|_2
 \quad(t\in[0,1],\ z,z'\in Z).
 \tag{150.1}
\]
Given any finite set of prescribed times in $[0,1]$, there is an
increasing definable homeomorphism $h:[0,1]\to[0,1]$ fixing those
times, and a finite $K\ge0$, such that
\[
 |F(h(s),z)-F(h(s'),z')|
 \le K|s-s'|+L\|z-z'\|_2.
 \tag{150.2}
\]
In particular the reparameterized family is uniformly Lipschitz in
its scalar parameter.

Consequently, if $A_t$ is a Hausdorff-continuous definable family of
nonempty compact sets contained in one fixed compact box in
$\mathbb R^n$, then such an $h$, fixing the same prescribed finite
times, makes
\[
 d_H(A_{h(s)},A_{h(s')})\le K|s-s'|.
 \tag{150.3}
\]
This consequence concerns the Hausdorff distance of the sets. It does
not construct homeomorphic point matching between them.
\end{saresult}

\begin{proof}
\textbf{Reduction to a full spatial box.}\quad
If $Z$ is empty, take $h=\operatorname{id}$; the function assertion
is vacuous. Otherwise enlarge the ambient dimension if necessary so
that $n\ge1$, and choose a full-dimensional compact box $Q$ containing
$Z$. Extend $F$ to $[0,1]\times Q$ by
\[
 \widetilde F(t,x)=
 \min_{z\in Z}\bigl(F(t,z)+L\|x-z\|_2\bigr).
 \tag{150.4}
\]
The minimum exists by compactness, and the formula is definable.
The triangle inequality shows that each spatial slice of
$\widetilde F$ is $L$-Lipschitz. Inequality (150.1) shows that it
agrees with $F$ on $Z$. Moreover,
\[
 |\widetilde F(t,x)-\widetilde F(t',x')|
 \le \sup_{z\in Z}|F(t,z)-F(t',z)|+L\|x-x'\|_2.
\]
Uniform continuity of $F$ on its compact domain proves joint
continuity of the extension. It is bounded on $[0,1]\times Q$.
We therefore work with this extension and write it as $F$.
At the end we restrict the conclusion to the original set $Z$.

\textbf{The exact original generic-Lipschitz input.}\quad
We use Guillaume Valette, \emph{Regular vectors and bi-Lipschitz
trivial stratifications in o-minimal structures},
arXiv:2110.05759v7, 29 September 2022, Proposition~5.7,
printed pp.~25--26. Its hypothesis is that each spatial function
in a definable family is Lipschitz. Its conclusion is a finite
partition of the parameter space on whose pieces the full function
is Lipschitz on every compact subset of the product with spatial
space. The introduction to Section~5, printed p.~22, explicitly
excepts this proposition from that section's polynomial-boundedness
restriction. We therefore use it in any fixed o-minimal expansion
of the real field, without assuming polynomial boundedness.

Applying that proposition and refining its one-dimensional parameter
partition, we obtain finitely many open intervals $J$ and isolated
times such that $F$ is Lipschitz on $J'\times Q$ whenever $J'$ is a
closed subinterval of $J$. Add the prescribed times and the endpoints
$0,1$ to this finite partition. No uniform Lipschitz bound up to the
endpoints of an open interval is being assumed.

\textbf{A definable bound for the time derivative.}\quad
On one of these open intervals $J$, choose an open dense definable
set $\Omega\subset J\times\operatorname{Int}Q$ on which $F$ is
$C^1$, by definable $C^1$ cell decomposition. Its complement in that
open product has dimension at most $n$. Set
\[
 M(t)=\sup\bigl(\{|\partial_t F(t,z)|:
                      (t,z)\in\Omega\}\cup\{0\}\bigr).
 \tag{150.5}
\]
This is a finite definable nonnegative function. Indeed, for any
$t\in J$ choose a closed interval $J'\subset J$ with $t$ in its
interior. Valette's conclusion gives a finite Lipschitz constant on
$J'\times Q$, bounding every derivative occurring in (150.5).
The supremum is definable by the ordered-field quantifiers.
The spatial bound (150.1) also gives
\[
 |\partial_{z_i}F(t,z)|\le L
 \quad((t,z)\in\Omega,\ 1\le i\le n).
 \tag{150.6}
\]

Partition $J$ further into finitely many points and intervals on
which either $M=0$ or $M>0$. On an interval where $M>0$, definable
choice gives $z(t)\in\operatorname{Int}Q$ with
\[
 (t,z(t))\in\Omega,\qquad
 |\partial_tF(t,z(t))|>\tfrac12 M(t).
 \tag{150.7}
\]
There is no assertion that this initial choice is continuous.
Put $u(t)=F(t,z(t))$. A further finite one-dimensional monotonicity
and $C^1$ partition makes $u$ and every coordinate $z_i$ both $C^1$
and monotone on each remaining open interval $(a,b)$. All these
functions are bounded, so they have finite one-sided limits at
both endpoints. By the chain rule and (150.6)--(150.7),
\[
 M(t)\le 2|u'(t)|+2L\sum_{i=1}^n|z_i'(t)|
 \qquad(a<t<b).
 \tag{150.8}
\]
Only the derivative at a full $C^1$ point is used in the chain rule.

\textbf{The clock has a finite endpoint without integration.}\quad
For such an interval set
\[
 \begin{split}
 \Phi(t)={}&t-a+2|u(t)-u(a+)|\\
           &+2L\sum_{i=1}^n|z_i(t)-z_i(a+)|.
 \end{split}
 \tag{150.9}
\]
Because all displayed functions are monotone, the derivatives of
the absolute differences are $|u'|$ and $|z_i'|$. Thus
\[
 \Phi'(t)\ge1+M(t)>0.
 \tag{150.10}
\]
The function extends continuously to $[a,b]$, with $\Phi(a)=0$
and with the finite endpoint
\[
 C=\Phi(b)=b-a+2|u(b-)-u(a+)|
              +2L\sum_i|z_i(b-)-z_i(a+)|>0.
 \tag{150.11}
\]
It is consequently an increasing definable homeomorphism
$[a,b]\to[0,C]$. On an interval where $M=0$, instead put
$\Phi(t)=t-a$ and $C=b-a$. In either case
$\Phi'\ge1+M$ throughout the open interval.
This explicit finite-variation formula does not use an integral
whose definability would have to be established.

\textbf{The derivative estimate gives uniform parameter control.}\quad
We justify this step at every spatial point, rather than assuming
that slice continuity already gives the needed estimate.
Let $E=(J\times\operatorname{Int}Q)\setminus\Omega$.
Since $\dim E\le n$, the definable set of $z$ for which
$E\cap(J\times\{z\})$ has positive dimension has dimension at
most $n-1$. Outside this exceptional spatial set a vertical line
meets $E$ in finitely many points. These good $z$ form a dense subset
of $\operatorname{Int}Q$.

Fix a good $z$. On each component of $(a,b)$ after its finitely many
exceptional times have been removed, the chain rule and (150.10)
show that, in the clock coordinate $w=\Phi(t)$,
\[
 \left|\frac{d}{dw}F(\Phi^{-1}(w),z)\right|
   =\frac{|\partial_tF(t,z)|}{\Phi'(t)}\le1.
 \tag{150.12}
\]
The ordinary mean value theorem gives the $1$-Lipschitz estimate
on each such component. Continuity across its finitely many endpoints
then gives
\[
 |F(t,z)-F(t',z)|\le|\Phi(t)-\Phi(t')|
 \quad(t,t'\in[a,b]).
 \tag{150.13}
\]
This includes the case $M=0$. Passing to limits from the dense set
of good spatial points, and then to the boundary of $Q$, proves
(150.13) for every $z\in Q$. Thus the bound is uniform in $z$.

Normalize this interval's clock to the original interval by
\[
 k_{a,b}(t)=a+\frac{b-a}{C}\Phi(t),\qquad
 h_{a,b}=k_{a,b}^{-1}.
 \tag{150.14}
\]
Both endpoints are fixed, and (150.13) gives
\[
 |F(h_{a,b}(s),z)-F(h_{a,b}(s'),z)|
 \le\frac{C}{b-a}|s-s'|
 \quad(s,s'\in[a,b],\ z\in Q).
 \tag{150.15}
\]
There are only finitely many final intervals. Glue their maps
$h_{a,b}$ at their shared endpoints to obtain an increasing
continuous definable bijection $h:[0,1]\to[0,1]$, hence a
homeomorphism. It fixes every partition point and in particular
every prescribed time. Let $K$ be the maximum of the finitely many
constants $C/(b-a)$. For two arbitrary parameter values, split the
interval between them at the finitely many partition points and
sum (150.15). The subinterval lengths sum to the original distance,
so
\[
 |F(h(s),z)-F(h(s'),z)|\le K|s-s'|.
\]
Combining this with the unchanged spatial bound proves (150.2),
and restriction to $Z$ proves the function assertion.

\textbf{The Hausdorff consequence (A150H).}\quad
Let all the nonempty compact fibres $A_t$ lie in a fixed compact box
$Q$. The distance function
\[
 d(t,z)=\operatorname{dist}(z,A_t),\qquad (t,z)\in[0,1]\times Q,
 \tag{150.16}
\]
is definable and spatially $1$-Lipschitz. Hausdorff continuity implies
joint continuity, because
\[
 |d(t,z)-d(t',z')|
 \le d_H(A_t,A_{t'})+\|z-z'\|_2.
\]
For any two nonempty compact subsets $A,B\subset Q$ one has the
exact identity
\[
 d_H(A,B)=\sup_{z\in Q}
       |\operatorname{dist}(z,A)-\operatorname{dist}(z,B)|.
 \tag{150.17}
\]
The inequality from right to left follows from the triangle inequality
for distance to a set. For the reverse inequality, each directed
Hausdorff distance is attained at a point of $A$ or $B$; at that point
one of the two distance functions is zero. The point lies in $Q$,
so the supremum is at least each directed distance.
Apply the function assertion to (150.16), and take the supremum in
$z$ using (150.17). This yields (150.3), with all prescribed times
still fixed.
\end{proof}

\paragraph{Original-source and scope record.}
The complete original source is
\url{https://arxiv.org/pdf/2110.05759}, explicitly version 7 dated
29 September 2022, 29 pages. Proposition~5.7 is on printed
pp.~25--26; Section~5's exception for this proposition is on printed
p.~22. The source was freshly opened and its full relevant statement
and proof read after recovery. This result concerns a single scalar
parameter over the real field. It does not assert a corresponding
homeomorphic reparameterization theorem for higher-dimensional
parameter cubes, and Hausdorff control alone does not produce a
homeomorphic trivialization.
